\documentclass[11pt]{article}
\usepackage[margin=1in]{geometry}
\usepackage{amsmath,amssymb}
\usepackage{microtype}
\usepackage[hidelinks]{hyperref}

\title{Why a Measure-Zero Periodic Set Can Contain Rational Parameters}
\author{}
\date{}

\begin{document}
\maketitle

Two elementary dynamical examples show why proving that periodic parameters are
rare, measure zero, or isolated does not by itself exclude rational parameters.

\section*{1. A measure-zero periodic set containing every rational}

Consider two uncoupled oscillators whose solution is
\[
  x(t)=\cos t,
  \qquad
  y(t)=\cos(ut),
  \qquad 0<u<1.
\]
The complete labelled state returns precisely when the two frequencies are
commensurable.  Equivalently,
\[
  u\in\mathbb{Q}.
\]
Thus the set of periodic parameters is
\[
  \mathbb{Q}\cap(0,1).
\]
This set is countable and has Lebesgue measure zero, yet it contains every
rational parameter in the interval.  It is also dense.  Consequently, a theorem
saying only that periodic parameters have measure zero would say nothing about
whether the Pythagorean parameters are periodic.

\section*{2. An isolated periodic parameter that is rational}

Consider
\[
  \ddot{x}+x=0,
  \qquad
  \ddot{y}+2y=0,
\]
with parameter-dependent initial conditions
\[
  x(0)=1,
  \quad \dot{x}(0)=0,
  \qquad
  y(0)=u-\frac13,
  \quad \dot{y}(0)=0.
\]
The resulting solution is
\[
  x(t)=\cos t,
  \qquad
  y(t)=\left(u-\frac13\right)\cos(\sqrt{2}\,t).
\]
For \(u\ne 1/3\), both incommensurable frequencies \(1\) and \(\sqrt{2}\) are
present, so the full orbit is quasiperiodic and never returns exactly.  At
\[
  u=\frac13,
\]
the second oscillator has zero amplitude, leaving the periodic motion
\(x(t)=\cos t\).  Hence the periodic-parameter set is exactly
\[
  \left\{\frac13\right\},
\]
an isolated, measure-zero set consisting of a rational number.

\section*{Relevance to the Pythagorean--Burrau conjecture}

Showing that intersections with the periodic-brake locus are isolated would be
substantial geometric progress, but it would not finish the conjecture.  One
would still need an exact bridge excluding rational intersection parameters---for
example,
\[
  \mathcal{B}(u,t)=0 \implies u\notin\mathbb{Q},
\]
or the stronger real statement
\[
  \mathcal{B}(u,t)\ne0
  \qquad
  \text{for every real }u\in(0,1)\text{ and every admissible }t>0.
\]

\end{document}
